\hypertarget{servicing-policies-with-estimated-action-and-observation-probabilities}{%
\chapter{Servicing policies with estimated action and observation
probabilities}\label{servicing-policies-with-estimated-action-and-observation-probabilities}}

\label{exec13:servicing-observed}

A servicing team wants to compare an assistance rule with ordinary
collection. Its historical action probabilities may not have been
retained, and some customers disappear from the analytical records
before their next repayment is measured. These are two estimation
problems only after a more basic scientific question has been answered:
did the recorded history include enough information to remove action and
observation selection? A highly accurate model of recorded actions
cannot answer that question on its own.

This chapter develops a restricted route with estimated probabilities, a
policy selected on separate training histories, and an explicit
finite-sample confidence statement. It also constructs two
observationally indistinguishable worlds in which that route answers
different causal questions. The confidence construction and finite
examples are project derivations. The longitudinal treatment/observation
identification principle is inherited from Hernán and Robins; the
importance of independent policy evaluation and cumulative, rather than
current-action, ratios is also explicit in Jiang and Li. Their results
do not themselves prove the composition below.
\citep{hernan2020,jiang2016,exec13_servicing_argument,exec13_servicing_exact}

\hypertarget{the-servicing-information-clock-and-the-eligible-population}{%
\section{The servicing information clock and the eligible
population}\label{the-servicing-information-clock-and-the-eligible-population}}

Fix a finite horizon \(T\). An independently sampled customer has a
complete potential history with pre-action information \(H_t\), binary
assistance \(A_t\), and subsequent net contribution \(R_t\), for
\(t=0,\ldots,T-1\). History contains every recorded past variable used
by the rule or the assignment mechanism, including earlier rewards when
relevant. A finite table, or an explicitly sufficient finite summary,
indexes the probability models. Arbitrary discretization of a continuous
history does not establish sufficiency.

Let \(O_{-1}=1\). Among records still observed at date \(t\), the bank
first records \(H_t\), then selects \(A_t\), then the observation
indicator \(C_t\) determines whether the following reward and next
history are recorded. Set \(O_t=O_{t-1}C_t\). Dropout is absorbing. Thus
\(C_t=1\) denotes retention here, the opposite coding from a
conventional censoring indicator. Define

\[
\begin{aligned}
 e_t(a\mid h)&=\Pr(A_t=a\mid H_t=h,O_{t-1}=1),\\
 g_t(h,a)&=\Pr(C_t=1\mid H_t=h,A_t=a,O_{t-1}=1).
\end{aligned}
\]

The stipulated observation law has four substantive restrictions. First,
consistency connects the recorded actions to the defined assistance
intervention. Second, conditional on \(H_t\), action selection is
independent of subsequent potential histories under the candidate
policies. Third, conditional on \(H_t,A_t\), retention is independent of
the subsequent complete history; dropout merely hides records and does
not change repayment. Fourth, all histories reachable under the
candidate policies satisfy known floors \(e_t(a\mid h)\ge c_e>0\) for
both actions and \(g_t(h,a)\ge c_g>0\), with \(c_e\le1/2\). Infeasible
actions must instead be removed from both the declared policy class and
the support requirement.

The independent sampling unit is the whole customer history. Dependence
within that history is allowed. Independent training and evaluation
cohorts follow the same complete-history, assignment and observation
law. Distinct customer identifiers do not guarantee this condition when
common calendar shocks, household links or shared staff information
remain. The construction does not qualify those sampling regimes.

The target is

\[
 V(\pi)=\mathbb E_\pi\left[\sum_{t=0}^{T-1} R_t\right],
 \qquad |R_t|\le B_t,
\]

per original eligible customer, with known finite \(B_t\). Discounting
and a bounded terminal account can be included in the dated rewards. A
sale, default or account closure that changes the economic outcome is
not automatically analytical dropout. Treating it as missing information
would change the target. The present example assigns neither a currency
value to customer welfare nor a bank-specific loss threshold.

\hypertarget{identification-with-action-and-retention-prefixes}{%
\section{Identification with action and retention
prefixes}\label{identification-with-action-and-retention-prefixes}}

For a history-based policy \(\pi_t(a\mid h)\), define

\[
 W_t^\pi=\prod_{s=0}^{t}
 \frac{\pi_s(A_s\mid H_s)}{e_s(A_s\mid H_s)}
 \frac{C_s}{g_s(H_s,A_s)}.
\]

An observed score stops at dropout; its subsequent contributions are
zero, without accessing unrecorded histories or evaluating undefined
denominators. Under the stated restrictions,

\[
 V(\pi)=\sum_{t=0}^{T-1}\mathbb E[W_t^\pi R_t].
\]

To verify the identity, factor a retained prefix into its
initial-history law, action probabilities, retention probabilities and
transition/reward kernels. Multiplication by the weight cancels each
action and retention factor and replaces each action factor with the
target policy. Sequential exchangeability and the observation
restriction make the remaining transition/reward kernels those of the
complete intervention law. Summing that prefix gives
\(\mathbb E_\pi R_t\). The observation factor for date \(t\) is required
because dropout occurs before that reward is recorded. Weighting only by
the current action leaves the incoming history distributed under
historical practice.

This argument uses the entire history law and does not require a Markov
state. A Markov reduction is an additional restriction that must be
checked in the selected construction. Nor does fitting \(e_t\) supply
the exchangeability premise: the observed conditional probability exists
even in a confounded world.

\hypertarget{a-finite-table-permits-an-explicit-probability-error-budget}{%
\section{A finite table permits an explicit probability error
budget}\label{a-finite-table-permits-an-explicit-probability-error-budget}}

Use \(m\) independent training histories to fit action and retention
probabilities. The same training set may select the policy; the
evaluation sample is independent of every such selection. Before
training, fix all \(J\) Bernoulli table cells: an action-one cell for
each relevant time/history and a retention cell for each relevant
time/history/action. Estimate each probability by its observed fraction,
conditional on reaching that cell. A customer contributes at most one
observation to a given dated cell.

Let \(N_j\) be that cell's training count and allocate total nuisance
failure probability \(\alpha_N\). For \(N_j>0\), put

\[
 \epsilon_j=\min\left\{1,\sqrt{\frac{\log(2J/\alpha_N)}{2N_j}}\right\}.
\]

For an empty cell choose any fixed estimate and set \(\epsilon_j=1\).
Conditional on the membership indicators for one cell, its observed
outcomes are independent Bernoulli variables with that cell's fixed
probability. Hoeffding's inequality gives failure probability at most
\(\alpha_N/J\); averaging over its random count preserves the bound.
Different cells can be dependent because customers visit several times:
a union bound, not cell independence, gives simultaneous error control
at least \(1-\alpha_N\).

For completeness, the concentration step follows from the
bounded-variable exponential-moment bound
\(\mathbb E e^{\lambda(X-\mathbb EX)}\le e^{\lambda^2/8}\) for
\(0\le X\le1\). Independence, exponential Markov inequality, and
optimization at \(\lambda=4x\) give \(\Pr(|\bar X-p|>x)\le2e^{-2Nx^2}\).
The bound follows by convexity of the exponential and bounding the
second derivative of the log Bernoulli moment-generating function by
\(1/4\).

Project each action-one estimate onto \([c_e,1-c_e]\) and obtain the
action-zero estimate by complement. Project retention onto \([c_g,1]\).
Projection cannot increase error when the true probability is in the
interval. It does not turn an unverified support floor into a scientific
fact. Let \(\epsilon_{e,t}\) and \(\epsilon_{g,t}\) be the largest
corresponding cell radii at time \(t\). On the simultaneous event,

\[
 \left|\frac{e_t}{\widehat e_t}-1\right|
 \le\frac{\epsilon_{e,t}}{c_e},\qquad
 \left|\frac{g_t}{\widehat g_t}-1\right|
 \le\frac{\epsilon_{g,t}}{c_g}.
\]

These statements hold for every policy measurable with respect to
training, including a rule selected using the same estimated
probabilities. There is no independence requirement between the policy
and those fits. Adaptive table construction, future-driven history
definitions and pooled time cells without a stationarity assumption are
excluded from this proof.

\hypertarget{propagating-fitted-probabilities-into-a-valid-interval}{%
\section{Propagating fitted probabilities into a valid
interval}\label{propagating-fitted-probabilities-into-a-valid-interval}}

For the frozen training-selected rule, use the fitted denominators in
the prefix score \(Y_i^\pi=\sum_t\widehat W_{it}^\pi R_{it}\), and
average it over \(n\) evaluation histories. Conditional on training,
cancellation of the true observed law gives

\[
 \mathbb E[Y_i^\pi\mid D_{\rm train}]
 =\sum_t\mathbb E_\pi\left[R_t
 \prod_{s=0}^t\frac{e_sg_s}{\widehat e_s\widehat g_s}\right].
\]

For nonnegative factors \(r_k\) with \(|r_k-1|\le d_k\), telescoping the
product gives \(|\prod_k r_k-1|\le\prod_k(1+d_k)-1\). This remains valid
when a lower bound \(1-d_k\) is negative. Therefore the conditional bias
is at most

\[
 D=\sum_t B_t\left[
 \prod_{s=0}^t
 \left(1+\frac{\epsilon_{e,s}}{c_e}\right)
 \left(1+\frac{\epsilon_{g,s}}{c_g}\right)-1\right].
\]

Independently of the nuisance event, each evaluation score lies in
\([-M,M]\), where

\[
 M=\sum_t B_t(c_ec_g)^{-(t+1)}.
\]

Applying the same bounded-variable argument after rescaling this range,
the radius

\[
 r_E=M\sqrt{\frac{2\log(2/\alpha_E)}{n}}
\]

bounds evaluation error with probability at least \(1-\alpha_E\),
conditional on any training sample. Consequently

\[
 I_\pi=[\bar Y^\pi-D-r_E,\ \bar Y^\pi+D+r_E]
 \cap[-B_\Sigma,B_\Sigma],\qquad B_\Sigma=\sum_t B_t,
\]

covers \(V(\pi_{D_{\rm train}})\) with unconditional probability at
least \(1-\alpha_N-\alpha_E\). This is a confidence statement about the
value of the realized trained policy. It is not the value of an unknown
optimal policy or the expectation of a policy-learning algorithm over
fresh training sets. Conditional coverage over evaluation is controlled
on the nuisance event; arbitrary bad training samples do not acquire a
conditional guarantee.

An empty intersection must be reported as an incompatible result, not
replaced by a point or silently dropped. For known action or retention
laws set their radii to zero and use their actual lower bounds. The
conservative total interval remains valid with estimated laws; an
evaluation-only normal standard error generally ignores the training
bias when both sample sizes are comparable.

For fixed \(J,T,c_e,c_g\), suppose each training cell has probability at
least \(p_{\min}>0\) of being reached. A multiplicative binomial bound
yields \(\Pr(\min_j N_j<mp_{\min}/2)\le J e^{-mp_{\min}/8}\). On its
complement, \(\max_j\epsilon_j\le\sqrt{\log(2J/\alpha_N)/(mp_{\min})}\).
Expanding the finite product shows
\(D=O(\sqrt{\log(J/\alpha_N)/(mp_{\min})})\), with constants depending
on horizon and support floors; \(r_E=O(n^{-1/2})\). Thus both
uncertainty terms can shrink in this finite class. Rare histories, a
growing horizon or a growing table can destroy useful precision. No
arbitrary-learner or continuous-history rate follows.

\hypertarget{policy-comparisons-preserve-the-selection-target}{%
\section{Policy comparisons preserve the selection
target}\label{policy-comparisons-preserve-the-selection-target}}

For \(K\) prespecified policies, or \(K\) rules constructed only from
training, replace \(\alpha_E\) by \(\alpha_E/K\). The nuisance event
already covers all their table cells. This gives simultaneous intervals
\([L_k,U_k]\), intersected with the value range, without asserting
independence between policy scores on the same customers.

The finite-class best value lies in \([\max_k L_k,\max_k U_k]\). If the
selected rule is member \(j\), its finite-class regret lies in

\[
 \left[\max\{0,\max_k L_k-U_j\},\ \max_k U_k-L_j\right].
\]

If any component interval is empty the family result is flagged
incompatible. Simultaneous coverage allows selecting a member after
inspecting these intervals, provided every candidate was fixed before
evaluation and all were covered. Adding candidates after viewing
evaluation outcomes needs a new argument. For an individual paired
contrast, one may instead bound the paired score and its nuisance error
directly; treating separately fitted values as independent loses their
covariance. Neither approach turns a finite candidate class into
unrestricted policy optimization.

\hypertarget{exact-small-references-and-the-cost-of-conservatism}{%
\section{Exact small references and the cost of
conservatism}\label{exact-small-references-and-the-cost-of-conservatism}}

A two-decision example uses observed stress \(S_t\in\{0,1\}\), with
initial probability \(1/2\). Historical assistance has probability
\(.3+.4S_t\). Next stress has probability \(.2+.5S_t-.15A_t\), retention
probability is \(.9-.1S_t-.1A_t\), and dated contribution is
\(2-3S_{t+1}-.25A_t\). These are hypothetical utility units; their
coefficient signs express an assistance cost and an improvement in
future stress. Every path probability is rational. Retention uses a draw
independent of the next stress innovation given current state and
action.

No assistance, always assist, and assist only when stressed form a fixed
finite policy class. Exhaustive path sums under interventions are
compared with a separate backward value recursion. Summation over
observed prefixes reproduces the same values with true denominators.
Perturbing action-one probabilities by \(+.04\) and retention
probabilities by \(-.03\) verifies the fitted-law bias expression and
its conservative envelope. A deliberately incorrect
current-action/retention shortcut is kept as a negative reference. The
program enumerates all paths; it draws no final evaluation sample and
measures no empirical coverage. \citep{exec13_servicing_exact}

\begin{longtable}[]{@{}lrr@{}}
\toprule
\begin{minipage}[b]{0.25\columnwidth}\raggedright
Policy\strut
\end{minipage} & \begin{minipage}[b]{0.33\columnwidth}\raggedleft
Exact intervention value\strut
\end{minipage} & \begin{minipage}[b]{0.33\columnwidth}\raggedleft
Bias from the specified fitted denominators\strut
\end{minipage}\tabularnewline
\midrule
\endhead
\begin{minipage}[t]{0.25\columnwidth}\raggedright
No assistance\strut
\end{minipage} & \begin{minipage}[t]{0.33\columnwidth}\raggedleft
1.3750\strut
\end{minipage} & \begin{minipage}[t]{0.33\columnwidth}\raggedleft
+.21997\strut
\end{minipage}\tabularnewline
\begin{minipage}[t]{0.25\columnwidth}\raggedright
Always assist\strut
\end{minipage} & \begin{minipage}[t]{0.33\columnwidth}\raggedleft
2.0000\strut
\end{minipage} & \begin{minipage}[t]{0.33\columnwidth}\raggedleft
-.22321\strut
\end{minipage}\tabularnewline
\begin{minipage}[t]{0.25\columnwidth}\raggedright
Assist when stressed\strut
\end{minipage} & \begin{minipage}[t]{0.33\columnwidth}\raggedleft
1.6625\strut
\end{minipage} & \begin{minipage}[t]{0.33\columnwidth}\raggedleft
+.20834\strut
\end{minipage}\tabularnewline
\bottomrule
\end{longtable}

The common worst-case bias envelope is 1.37243 utility units, exceeding
every realized bias in the table. That conservatism is the price of
using probability error bounds without favorable outcome cancellation.
These perturbations illustrate the error propagation formula; they are
not estimates learned from an executed training sample.

Even with \(B_t=2\), \(c_e=.25\), \(c_g=.65\), \(T=2\), \(n=1000\) and
\(\alpha_E=.025\), the evaluation radius alone exceeds eight utility
units. The entire stipulated value range is only eight units wide. A
mathematically valid bound can therefore be unsuitable for choosing
assistance. Increasing the nominal confidence level or adding more rare
history cells makes this tradeoff harder. This is a reason to seek
shorter questions, stronger information, sharper variance-sensitive
bounds or a reviewed sequential regression construction; it is not a
reason to remove the error budget.

Sequential doubly robust scores offer a substantive alternative when
continuation values are useful. They require a separate analysis of
fitted action/retention models, value-function errors and sampling
dependence. The known-logger variance reduction in Jiang and Li does not
automatically establish estimated-logger finite coverage. The earlier
synthetic collection results already show that mean protection and
normal-interval reliability differ. \citep{jiang2016}

\hypertarget{observational-equivalence-prevents-a-hidden-state-repair-by-prediction}{%
\section{Observational equivalence prevents a hidden-state repair by
prediction}\label{observational-equivalence-prevents-a-hidden-state-repair-by-prediction}}

Consider a single decision with no censoring. In world E, action is a
fair independent coin and \(Y=A\) with probability \(.8\), otherwise
\(Y=1-A\). The intervention values are \(V_E(1)=.8\), \(V_E(0)=.2\). In
world H, private hardship \(U\) is a fair coin, \(Y=U\), and
\(\Pr(A=1\mid U)=.2+.6U\). Here \(V_H(1)=V_H(0)=.5\).

Both worlds have exactly the same observed joint probabilities:

\begin{longtable}[]{@{}lrrrr@{}}
\toprule
Observed pair \((A,Y)\) & \((0,0)\) & \((0,1)\) & \((1,0)\) &
\((1,1)\)\tabularnewline
\midrule
\endhead
Either world & .4 & .1 & .1 & .4\tabularnewline
\bottomrule
\end{longtable}

Even the private-state assignment in world H has overlap: its
probabilities are .2 and .8. Infinite recorded data estimate
\(\Pr(A=1)=.5\) and \(\mathbb E[Y\mid A=1]=.8\) perfectly in both
worlds. Weighting by that observed probability returns .8, while the
hidden world's value is .5. No fitted logger, extra evaluation data or
language-model interpretation of these same records can distinguish the
worlds. The effect is .6 in one and zero in the other.

Observation selection has a separate obstruction. Suppose half the
cohort is observed and half of those observed have \(Y=1\). One
full-data world assigns zero to every missing outcome; another assigns
one. Both yield the same observed records and retention probability .5,
but their cohort means are .25 and .75. Accurately estimating the
retention probability does not recover the unobserved outcome law. The
complete-case mean .5 answers neither world's population question.

An additional randomized action design, a truthful relevant-state audit,
or a validated proxy/instrument model can distinguish some such worlds.
Each adds information or restrictions. The existing proximal and
partial-observation treatments explain specific routes and their rank,
exclusion and policy-information requirements; they do not authorize
arbitrary hidden-state inference.
\citep{bennettprox2023,tennenholtz2020}

\hypertarget{what-the-proposal-can-recommend-after-this-derivation}{%
\section{What the proposal can recommend after this
derivation}\label{what-the-proposal-can-recommend-after-this-derivation}}

The preserved selected17 study used truthful representative state audits
and supplied action probabilities. It passed its selected valid cells,
while wrong probabilities covered only one of 512 requests and policy
leakage returned no finite outputs. The present finite-table route makes
estimated-probability uncertainty explicit after restricting to
sufficient observed histories. It removes no hidden-state requirement
from that predecessor and claims no new repeated-sampling qualification.
\citep{exec13_servicing_predecessor}

The separate proposed continuous-state servicing experiment has Gaussian
innovations and a quadratic reward. Its rewards are unbounded and its
history table is not finite. Consequently it cannot use this
bounded-table theorem as written. A prospectively declared bounded
utility would define a different target; alternatively, a valid
tail/moment bound and a continuous-model estimation argument are needed.
Quietly capping realized losses after evaluating the results would
answer a changed question.

For the proposal, the executable next comparison is a declared
finite-state generated demonstration with known evaluator truth,
estimated training probabilities, independent evaluation and adverse
hidden/censoring cases. Its useful-width screen must be chosen in
utility units before final draws. A broad interval or correct refusal is
an informative demonstration outcome. A claim about real servicing
remains conditional on an information audit and a justified dependence
law, and an unknown optimal policy remains a separate research question.
These distinctions make a prototype experiment worth designing without
pretending that a bank trial or a universal hidden-state theorem has
already been supplied. The following temporal analysis asks a related question at the level of causal structure: which intervention can distinguish effects that the observed law leaves equivalent? The servicing examples show why that information question must precede a more elaborate probability fit.

\label{exec13:servicing-observed-end}
